\documentclass{article}
\usepackage{algorithm}
\usepackage{algorithmic}
\usepackage{amsmath}
\usepackage{listings}
\usepackage{xcolor}

\title{OpenSHMEM Message Aggregation Pass for MLIR}
\author{Michael Beebe}
\date{\today}

\begin{document}

\maketitle

\section{Introduction}

The OpenSHMEM Message Aggregation Pass is a compiler optimization implemented within the Multi-Level Intermediate Representation (MLIR) framework that analyzes and optimizes OpenSHMEM Remote Memory Access (RMA) operations. The pass identifies opportunities to coalesce multiple small communication operations into fewer, larger operations, thereby reducing communication overhead and improving application performance.

\section{Pass Architecture}

The message aggregation pass employs a multi-phase analysis and transformation approach consisting of four main components:

\subsection{Memory Access Analysis}
The \texttt{MemoryAccessAnalyzer} component extracts memory access patterns from OpenSHMEM operations, creating a unified representation that captures:
\begin{itemize}
\item Operation type (PUT, GET) and variant (MEMORY, TYPED, SIZED, POINT)
\item Target processing element (PE) and context information
\item Memory references for source and destination
\item Transfer sizes and blocking behavior
\item Element sizes for typed operations
\end{itemize}

\subsection{Memory Layout Analysis}
The \texttt{MemoryLayoutAnalyzer} performs sophisticated memory region analysis to determine:
\begin{itemize}
\item Contiguity of memory accesses through offset calculation
\item Stride patterns for regular non-contiguous accesses
\item Base address extraction from complex memory reference chains
\item Support for \texttt{memref.subview} and \texttt{memref.view} operations
\end{itemize}

\subsection{Communication Pattern Detection}
The \texttt{CommunicationPatternDetector} groups operations based on compatibility criteria:
\begin{itemize}
\item Same target PE and context
\item Identical blocking behavior (blocking vs. non-blocking)
\item Compatible operation types and variants
\item Matching element sizes for sized operations
\end{itemize}

\subsection{Transformation Engine}
The \texttt{TransformationEngine} applies coalescing optimizations using multiple strategies:
\begin{enumerate}
\item \textbf{Duplicate Removal}: Eliminates identical operations
\item \textbf{Contiguous Coalescing}: Merges operations accessing adjacent memory regions
\item \textbf{Non-blocking Batching}: Optimizes scheduling of non-blocking operations
\end{enumerate}

\section{Algorithm}

\begin{algorithm}
\caption{OpenSHMEM Message Aggregation}
\label{alg:message_aggregation}
\begin{algorithmic}[1]
\REQUIRE Function $F$ containing OpenSHMEM operations
\ENSURE Optimized function with coalesced operations

\STATE $accesses \leftarrow \emptyset$
\STATE $groups \leftarrow \emptyset$

\COMMENT{Phase 1: Extract Memory Access Patterns}
\FOR{each operation $op$ in $F$}
    \IF{$op$ is OpenSHMEM RMA operation}
        \STATE $access \leftarrow$ \textsc{ExtractMemoryAccess}($op$)
        \STATE $accesses \leftarrow accesses \cup \{access\}$
    \ENDIF
\ENDFOR

\COMMENT{Phase 2: Group Compatible Operations}
\FOR{each $access$ in $accesses$}
    \STATE $key \leftarrow (access.pe, access.ctx, access.isNonBlocking,$
    \STATE \hspace{2em} $access.opType, access.opVariant, access.elementSize)$
    \IF{$groups[key]$ exists}
        \STATE $groups[key].operations \leftarrow groups[key].operations \cup \{access.op\}$
    \ELSE
        \STATE $groups[key] \leftarrow$ \textsc{CreateGroup}($access$)
    \ENDIF
\ENDFOR

\COMMENT{Phase 3: Apply Coalescing Transformations}
\FOR{each $group$ in $groups$ where $|group.operations| \geq 2$}
    \IF{\textsc{AreAllOperationsIdentical}($group$)}
        \STATE \textsc{RemoveDuplicates}($group$)
    \ELSIF{\textsc{CanCreateCoalescedOperation}($group$)}
        \STATE \textsc{CreateCoalescedOperation}($group$)
    \ELSIF{$group.isNonBlocking$ \AND \textsc{CanBatch}($group$)}
        \STATE \textsc{BatchOperations}($group$)
    \ENDIF
\ENDFOR

\RETURN optimized function $F'$
\end{algorithmic}
\end{algorithm}

\section{Implementation Details}

\subsection{Operation Coverage}
The pass supports all major OpenSHMEM RMA operations:
\begin{itemize}
\item \textbf{Memory operations}: \texttt{putmem}, \texttt{getmem}, \texttt{putmem\_nbi}, \texttt{getmem\_nbi}
\item \textbf{Typed operations}: \texttt{put}, \texttt{get}, \texttt{put\_nbi}, \texttt{get\_nbi}
\item \textbf{Context-aware variants}: \texttt{ctx\_put}, \texttt{ctx\_get}, \texttt{ctx\_put\_nbi}, \texttt{ctx\_get\_nbi}
\item \textbf{Sized operations}: \texttt{put8/16/32/64/128}, \texttt{get8/16/32/64/128}
\item \textbf{Point operations}: \texttt{p}, \texttt{g}
\end{itemize}

\subsection{Safety Constraints}
The transformation engine ensures correctness through several safety mechanisms:
\begin{enumerate}
\item \textbf{Memory Safety}: Operations are only coalesced if they access compatible memory regions
\item \textbf{Ordering Preservation}: OpenSHMEM ordering semantics within contexts are maintained
\item \textbf{Type Compatibility}: Only operations with identical types and variants are combined
\item \textbf{Context Isolation}: Operations using different contexts are never coalesced
\end{enumerate}

\subsection{Configuration Options}
The pass provides several configuration parameters:
\begin{itemize}
\item \texttt{enablePutCoalescing}: Enable/disable PUT operation coalescing
\item \texttt{enableGetCoalescing}: Enable/disable GET operation coalescing  
\item \texttt{minMessageSize}: Minimum message size threshold for coalescing
\item \texttt{maxCoalescingDistance}: Maximum distance between coalescable operations
\item \texttt{debugMode}: Enable detailed debug output
\end{itemize}

\section{Experimental Results}

The pass was evaluated on a comprehensive test suite of 14 test cases covering various OpenSHMEM operation patterns. Key results include:

\begin{itemize}
\item \textbf{Duplicate Elimination}: Successfully removes redundant identical operations
\item \textbf{Type Safety}: Correctly prevents coalescing of incompatible operation types (e.g., \texttt{put32} vs \texttt{put64})
\item \textbf{Context Awareness}: Properly handles context-based operation separation
\item \textbf{Blocking Behavior}: Maintains distinction between blocking and non-blocking operations
\end{itemize}

\section{Future Work}

The current implementation provides a solid foundation for advanced optimizations:

\begin{enumerate}
\item \textbf{Contiguous Memory Coalescing}: Merge operations accessing adjacent memory regions into single larger transfers
\item \textbf{Cross-operation Type Optimization}: Combine multiple small operations into larger sized operations (e.g., \texttt{put32} + \texttt{put32} $\rightarrow$ \texttt{put64})
\item \textbf{Synchronization Boundary Analysis}: Optimize across synchronization points while preserving correctness
\item \textbf{Performance Heuristics}: Implement cost models to guide transformation decisions
\end{enumerate}

\section{Conclusion}

The OpenSHMEM Message Aggregation Pass demonstrates the effectiveness of MLIR's pattern rewriting infrastructure for domain-specific optimizations. By providing a comprehensive framework for analyzing and optimizing OpenSHMEM communication patterns, the pass enables significant performance improvements while maintaining correctness guarantees. The modular architecture and extensive test coverage provide a robust foundation for future research in HPC compiler optimizations.

\end{document}
